\chapter{Topology}


Topology is often described as geometry without measurement. This slogan is useful but incomplete. A better description is that topology studies which properties of spaces are preserved by continuous deformation. Lengths, angles, and areas may change, but nearness, separation, connectedness, compactness, and continuity remain meaningful.

The transition from metric spaces to topological spaces is motivated by abstraction. In a metric space, open sets are defined using distances. But many arguments in analysis use only the open sets themselves, not the numerical values of the distances. Topology isolates this open-set structure and asks what can still be proved once distance is removed from the foundations.

Topology is the study of the qualitative structure of spaces. In Euclidean geometry, we measure length, angle, and area. In metric spaces, we measure distance. In topology, we keep only the information needed to speak about nearness, convergence, continuity, connectedness, and compactness.

This may seem like a loss of information, but it is in fact a powerful abstraction. Many theorems in analysis do not depend on the exact numerical value of a distance. The Intermediate Value Theorem depends on connectedness. The Extreme Value Theorem depends on compactness. Continuity can be described without mentioning a formula for distance. Topology isolates these structural ideas.

The basic object is a set together with a chosen collection of subsets called open sets. Once the open sets are specified, the language of limits and continuity can be rebuilt.

\section{From Metric Spaces to Topological Spaces}

\subsection{Open Sets in Metric Spaces}


In a metric space, the open ball is the primitive local object. Saying that a set is open means that every point of the set has a small protective neighborhood still contained in the set. This captures the intuitive idea that points inside an open set are not sitting on the edge. Around each point there is room to move slightly without leaving the set.

Notice that openness is relative to the ambient space. The interval $(0,1)$ is open in $\mathbb{R}$, while the same set viewed inside $[0,1]$ is also open in the subspace topology because it equals $[0,1]\cap (0,1)$. This dependence on context is not a defect; it is exactly what allows topology to describe spaces internally.

Let us recall the metric-space intuition. If $(X,d)$ is a metric space, then an open ball centered at $x\in X$ with radius $r>0$ is
\[
    B_r(x)=\{y\in X:d(x,y)<r\}.
\]
A subset $U\subseteq X$ is open if for every $x\in U$, there exists $r>0$ such that $B_r(x)\subseteq U$.

\begin{example}
In $\mathbb{R}$ with the usual metric $d(x,y)=|x-y|$, the interval $(0,1)$ is open. For every $x\in(0,1)$, we may choose
\[
    r<\min\{x,1-x\},
\]
so that $(x-r,x+r)\subseteq(0,1)$.
\end{example}

The important observation is that once we know which subsets are open, we can define continuity without explicitly using the metric.

\begin{proposition}[Topological Form of Continuity in Metric Spaces]
Let $(X,d_X)$ and $(Y,d_Y)$ be metric spaces. A function $f:X\to Y$ is continuous if and only if for every open set $V\subseteq Y$, the preimage $f^{-1}(V)$ is open in $X$.
\end{proposition}

\begin{proof}
Assume first that $f$ is $\varepsilon$-$\delta$ continuous. Let $V\subseteq Y$ be open and let $x\in f^{-1}(V)$. Then $f(x)\in V$, so there exists $\varepsilon>0$ such that $B_\varepsilon(f(x))\subseteq V$. By continuity, there exists $\delta>0$ such that $d_X(x,x')<\delta$ implies $d_Y(f(x),f(x'))<\varepsilon$. Thus $B_\delta(x)\subseteq f^{-1}(V)$, so $f^{-1}(V)$ is open.

Conversely, assume preimages of open sets are open. For $x\in X$ and $\varepsilon>0$, the ball $B_\varepsilon(f(x))$ is open in $Y$, so $f^{-1}(B_\varepsilon(f(x)))$ is open in $X$ and contains $x$. Therefore it contains some ball $B_\delta(x)$, giving the usual $\varepsilon$-$\delta$ condition.
\end{proof}

This proposition motivates the axiomatic definition of topology.

\section{Topological Spaces}

\subsection{Definition of a Topology}


A topology keeps only the information needed to speak about openness. The axioms say that the whole space and the empty set are open, that arbitrary unions of open sets are open, and that finite intersections of open sets are open. These rules match the behavior of open sets in metric spaces, but they no longer require a distance function.

The asymmetry between arbitrary unions and finite intersections is important. Arbitrary unions of open sets remain open because a point belongs to at least one open set and can use that set's local neighborhood. Infinite intersections may fail to be open because the available neighborhoods may shrink to zero size. For instance, $\bigcap_{n=1}^{\infty}(-1/n,1/n)=\{0\}$ is not open in $\mathbb{R}$.

\begin{definition}[Topology]
Let $X$ be a set. A \textbf{topology} on $X$ is a collection $\mathcal{T}\subseteq\mathcal{P}(X)$ satisfying:
\begin{enumerate}
    \item $\emptyset\in\mathcal{T}$ and $X\in\mathcal{T}$;
    \item arbitrary unions of elements of $\mathcal{T}$ belong to $\mathcal{T}$;
    \item finite intersections of elements of $\mathcal{T}$ belong to $\mathcal{T}$.
\end{enumerate}
The pair $(X,\mathcal{T})$ is called a \textbf{topological space}. Elements of $\mathcal{T}$ are called \textbf{open sets}.
\end{definition}

The asymmetry between arbitrary unions and finite intersections is essential. In $\mathbb{R}$, arbitrary unions of open intervals are open, but arbitrary intersections of open sets need not be open.

\begin{example}
For each $n\in\mathbb{N}$, the interval $(-1/n,1/n)$ is open in $\mathbb{R}$. However,
\[
    \bigcap_{n=1}^{\infty}(-1/n,1/n)=\{0\},
\]
and $\{0\}$ is not open in the usual topology of $\mathbb{R}$.
\end{example}

\subsection{Fundamental Examples}


The examples in this subsection should not be memorized as a disconnected list. They show that a topology can be coarse or fine depending on how many open sets are allowed. A coarse topology has fewer open sets, so it distinguishes fewer local behaviors. A fine topology has more open sets, so it can distinguish more local behaviors. The indiscrete topology is maximally coarse, while the discrete topology is maximally fine.

This comparison also explains why continuity depends on the chosen topology. If the domain topology is very fine, many preimages are open and continuity becomes easier. If the codomain topology is very fine, there are more open sets whose preimages must be checked and continuity becomes harder. Therefore topology is not just a property of a set; it controls which functions out of and into that set are considered continuous.

\begin{example}[Discrete Topology]
For any set $X$, the collection $\mathcal{T}=\mathcal{P}(X)$ is a topology. Every subset is open. This is called the \textbf{discrete topology}.
\end{example}

\begin{example}[Indiscrete Topology]
For any set $X$, the collection $\mathcal{T}=\{\emptyset,X\}$ is a topology. This is called the \textbf{indiscrete topology} or \textbf{trivial topology}.
\end{example}

\begin{example}[Usual Topology on $\mathbb{R}^n$]
The topology generated by Euclidean open balls in $\mathbb{R}^n$ is called the usual topology. Unless otherwise stated, $\mathbb{R}^n$ is equipped with this topology.
\end{example}

\begin{example}[Cofinite Topology]
Let $X$ be a set. Define
\[
    \mathcal{T}=\{\emptyset\}\cup\{U\subseteq X:X\setminus U\text{ is finite}\}.
\]
Then $\mathcal{T}$ is a topology on $X$, called the \textbf{cofinite topology}.
\end{example}

\begin{proof}
The empty set belongs to $\mathcal{T}$, and $X\setminus X=\emptyset$ is finite, so $X\in\mathcal{T}$. If $\{U_i\}_{i\in I}$ are non-empty open sets, then
\[
    X\setminus\bigcup_{i\in I}U_i=\bigcap_{i\in I}(X\setminus U_i),
\]
which is a subset of one finite set $X\setminus U_j$ if at least one $U_j$ is non-empty; hence it is finite. Thus arbitrary unions are open. If $U_1,\dots,U_n$ are open, then
\[
    X\setminus\bigcap_{k=1}^nU_k=\bigcup_{k=1}^n(X\setminus U_k),
\]
a finite union of finite sets, hence finite. Thus finite intersections are open.
\end{proof}

\begin{example}[Particular Point Topology]
Fix $p\in X$. Define
\[
    \mathcal{T}=\{\emptyset\}\cup\{U\subseteq X:p\in U\}.
\]
Then $\mathcal{T}$ is a topology. In this topology, every non-empty open set must contain the special point $p$.
\end{example}

\subsection{Comparing Topologies}


Given two topologies on the same set, saying that one is finer than the other means that it contains more open sets. This order relation is a way of comparing how much local information the topologies remember. The finer topology can make distinctions that the coarser topology cannot.

One should be careful with the direction of continuity under this comparison. The identity map from a finer topology to a coarser topology is continuous, because every open set in the coarser topology is already open in the finer topology. The reverse identity map may fail to be continuous. This small observation is often the quickest way to check whether two topologies are genuinely different.

\begin{definition}[Finer and Coarser Topologies]
Let $\mathcal{T}_1$ and $\mathcal{T}_2$ be topologies on the same set $X$. We say $\mathcal{T}_1$ is \textbf{finer} than $\mathcal{T}_2$ if
\[
    \mathcal{T}_2\subseteq\mathcal{T}_1.
\]
Equivalently, $\mathcal{T}_2$ is \textbf{coarser} than $\mathcal{T}_1$.
\end{definition}

A finer topology has more open sets. More open sets make it easier for subsets to be open, but harder for functions into the space to be continuous.

\begin{example}
On any set $X$, the discrete topology is the finest topology, while the indiscrete topology is the coarsest topology.
\end{example}

\begin{remark}
Topology is not only about sets; it is about sets together with a chosen standard of local visibility. A subset is open if the topology allows us to see it as a neighborhood-like object.
\end{remark}

\section{Bases and Subbases}


A basis is a compressed description of a topology. Instead of listing all open sets, we list a collection of basic open sets and then allow arbitrary unions. This mirrors the way intervals generate the usual topology on $\mathbb{R}$. Most open sets are too numerous and irregular to list directly, but they can be assembled from simple pieces.

The value of a basis is both conceptual and practical. Conceptually, it tells us what local neighborhoods look like. Practically, it allows us to prove continuity and openness by checking only basic open sets. In many spaces, especially product spaces, the topology is easiest to understand through its basis rather than through all open sets at once.

It is often inconvenient to list every open set. Instead, we describe a smaller collection of basic open sets from which all open sets can be built.

\subsection{Basis for a Topology}

\begin{definition}[Basis]
Let $X$ be a set. A collection $\mathcal{B}\subseteq\mathcal{P}(X)$ is a \textbf{basis} for a topology on $X$ if:
\begin{enumerate}
    \item for every $x\in X$, there exists $B\in\mathcal{B}$ such that $x\in B$;
    \item if $x\in B_1\cap B_2$ with $B_1,B_2\in\mathcal{B}$, then there exists $B_3\in\mathcal{B}$ such that
    \[
        x\in B_3\subseteq B_1\cap B_2.
    \]
\end{enumerate}
The topology generated by $\mathcal{B}$ consists of all unions of elements of $\mathcal{B}$.
\end{definition}

\begin{example}
Open intervals $(a,b)$ form a basis for the usual topology on $\mathbb{R}$. Open balls form a basis for the usual topology on $\mathbb{R}^n$.
\end{example}

\begin{proposition}
If $\mathcal{B}$ is a basis on $X$, then the collection of all unions of basis elements is a topology on $X$.
\end{proposition}

\begin{proof}
The empty union gives $\emptyset$, and the first basis condition ensures that $X$ is the union of all basis elements. Arbitrary unions of unions of basis elements are still unions of basis elements. For finite intersections, it is enough to consider two open sets $U$ and $V$. If $x\in U\cap V$, then $x\in B_1\subseteq U$ and $x\in B_2\subseteq V$ for some basis elements. By the second basis condition, there exists $B_3$ such that $x\in B_3\subseteq B_1\cap B_2\subseteq U\cap V$. Thus $U\cap V$ is a union of basis elements.
\end{proof}

\subsection{Subbasis}


A subbasis is an even more economical generating device than a basis. Starting from a subbasis, we first take finite intersections to form a basis, and then arbitrary unions to form the topology. The two-step procedure is important: finite intersections create local compatibility among the chosen subbasic conditions, while arbitrary unions allow open sets to be assembled from those local pieces.

Subbases are especially useful when a topology is defined by requiring many maps to be continuous. For example, product topologies can be described using inverse images of open sets under coordinate projection maps. The subbasis captures the minimum open-set information needed to make those projections continuous.

\begin{definition}[Subbasis]
A collection $\mathcal{S}\subseteq\mathcal{P}(X)$ is a \textbf{subbasis} for a topology on $X$ if the set of all finite intersections of elements of $\mathcal{S}$ forms a basis. The topology generated by $\mathcal{S}$ is the collection of arbitrary unions of such finite intersections.
\end{definition}

\begin{example}
In $\mathbb{R}$, the collection of all rays $(-\infty,a)$ and $(a,\infty)$ is a subbasis for the usual topology.
\end{example}

\subsection{Countability Axioms}

\begin{definition}[First and Second Countability]
A topological space $X$ is \textbf{first countable} if every point has a countable local basis. It is \textbf{second countable} if the topology has a countable basis.
\end{definition}

\begin{example}
The usual topology on $\mathbb{R}^n$ is second countable. A countable basis is given by all open balls whose centers have rational coordinates and whose radii are positive rational numbers.
\end{example}

\begin{proposition}
Every second countable space is first countable.
\end{proposition}

\begin{proof}
Let $\mathcal{B}$ be a countable basis for $X$. For a fixed point $x\in X$, the collection of all basis elements containing $x$ is countable. It is a local basis at $x$.
\end{proof}

\section{Closed Sets, Closure, Interior, and Boundary}


Closure, interior, and boundary are three ways of describing how a set sits inside a space. The interior consists of points safely inside the set. The closure consists of points that cannot be separated from the set by an open neighborhood. The boundary consists of points where every neighborhood sees both the set and its complement.

These notions are topological because they depend only on open sets. A point is in the closure of $A$ if every open neighborhood of the point intersects $A$. This definition does not mention distance, although in metric spaces it agrees with the more familiar sequential intuition: points in the closure can often be approached by sequences from the set.

Open sets are primary in the axioms, but closed sets are equally important.

\begin{definition}[Closed Set]
Let $X$ be a topological space. A subset $C\subseteq X$ is \textbf{closed} if $X\setminus C$ is open.
\end{definition}

\begin{proposition}[Properties of Closed Sets]
In a topological space $X$:
\begin{enumerate}
    \item $\emptyset$ and $X$ are closed;
    \item arbitrary intersections of closed sets are closed;
    \item finite unions of closed sets are closed.
\end{enumerate}
\end{proposition}

\begin{proof}
These properties follow from the axioms for open sets and De Morgan's laws:
\[
    X\setminus\bigcap_i C_i=\bigcup_i(X\setminus C_i),
    \qquad
    X\setminus\bigcup_{k=1}^n C_k=\bigcap_{k=1}^n(X\setminus C_k).
\]
\end{proof}

\subsection{Closure and Interior}

\begin{definition}[Closure]
Let $A\subseteq X$. The \textbf{closure} of $A$, denoted $\overline{A}$, is the intersection of all closed sets containing $A$.
\end{definition}

\begin{definition}[Interior]
Let $A\subseteq X$. The \textbf{interior} of $A$, denoted $A^\circ$ or $\operatorname{Int}(A)$, is the union of all open sets contained in $A$.
\end{definition}

\begin{definition}[Boundary]
The \textbf{boundary} of $A$ is
\[
    \partial A=\overline{A}\setminus A^\circ.
\]
Equivalently,
\[
    \partial A=\overline{A}\cap\overline{X\setminus A}.
\]
\end{definition}

\newpage

\begin{example}
In $\mathbb{R}$ with the usual topology, for $A=(0,1)$,
\[
    \overline{A}=[0,1],
    \qquad
    A^\circ=(0,1),
    \qquad
    \partial A=\{0,1\}.
\]
For $A=\mathbb{Q}\subseteq\mathbb{R}$,
\[
    \overline{\mathbb{Q}}=\mathbb{R},
    \qquad
    \operatorname{Int}(\mathbb{Q})=\emptyset.
\]
\end{example}

\begin{proposition}[Closure via Open Neighborhoods]
Let $A\subseteq X$ and $x\in X$. Then $x\in\overline{A}$ if and only if every open set $U$ containing $x$ intersects $A$.
\end{proposition}

\begin{proof}
If some open neighborhood $U$ of $x$ is disjoint from $A$, then $X\setminus U$ is a closed set containing $A$ but not containing $x$, so $x\notin\overline{A}$. Conversely, if $x\notin\overline{A}$, then $X\setminus\overline{A}$ is an open neighborhood of $x$ disjoint from $A$.
\end{proof}

\subsection{Limit Points and Density}

\begin{definition}[Limit Point]
Let $A\subseteq X$. A point $x\in X$ is a \textbf{limit point} of $A$ if every open neighborhood of $x$ contains a point of $A$ different from $x$.
\end{definition}

\begin{definition}[Dense Subset]
A subset $A\subseteq X$ is \textbf{dense} in $X$ if $\overline{A}=X$.
\end{definition}

\begin{example}
Both $\mathbb{Q}$ and $\mathbb{R}\setminus\mathbb{Q}$ are dense in $\mathbb{R}$ with the usual topology. Every open interval contains both rational and irrational numbers.
\end{example}

\section{Constructing New Spaces}


New spaces are often built from old ones. The subspace topology describes what openness means inside a subset. The product topology describes what openness means when several coordinates vary independently. The quotient topology describes what openness means after gluing points together. These three constructions are the basic grammar of topology.

It is helpful to associate each construction with a geometric operation. Subspaces correspond to cutting out a part of a space. Products correspond to forming coordinate systems or configuration spaces. Quotients correspond to identification, such as gluing the ends of an interval to form a circle. Many important spaces are most naturally defined by quotient constructions.

Topology becomes powerful when we can construct new spaces from old spaces.

\subsection{Subspace Topology}

\begin{definition}[Subspace Topology]
Let $(X,\mathcal{T})$ be a topological space and let $Y\subseteq X$. The \textbf{subspace topology} on $Y$ is
\[
    \mathcal{T}_Y=\{Y\cap U:U\in\mathcal{T}\}.
\]
\end{definition}

\begin{example}
The interval $[0,1]$ as a subspace of $\mathbb{R}$ has open sets of the form $[0,1]\cap U$, where $U$ is open in $\mathbb{R}$. For example, $[0,1/2)$ is open in $[0,1]$, although it is not open in $\mathbb{R}$.
\end{example}

\subsection{Product Topology}


The product topology is designed so that coordinate projections are continuous and so that open neighborhoods describe independent variation in finitely many coordinates. In a finite product such as $X\times Y$, basic open sets look like $U\times V$. In infinite products, however, basic open sets restrict only finitely many coordinates and leave all other coordinates unrestricted.

This finite restriction rule is not arbitrary. It is what makes compactness behave well under products and what prevents the topology from becoming too fine. The product topology therefore balances two needs: it remembers the topology of each coordinate while still allowing infinite products to retain useful global properties.

\begin{definition}[Product Topology]
Let $X$ and $Y$ be topological spaces. The \textbf{product topology} on $X\times Y$ is generated by the basis
\[
    \mathcal{B}=\{U\times V:U\subseteq X\text{ open},\ V\subseteq Y\text{ open}\}.
\]
\end{definition}

\begin{example}
The usual topology on $\mathbb{R}^2$ is the same as the product topology on $\mathbb{R}\times\mathbb{R}$, where each copy of $\mathbb{R}$ has the usual topology.
\end{example}

\begin{proposition}[Continuity into a Product]
A map $f:Z\to X\times Y$ is continuous if and only if its coordinate maps $\pi_X\circ f:Z\to X$ and $\pi_Y\circ f:Z\to Y$ are continuous.
\end{proposition}

\subsection{Quotient Topology}


The quotient topology is usually the least intuitive construction at first because open sets are defined indirectly. A set in the quotient is declared open exactly when its full preimage upstairs is open. This is the correct definition because the quotient map is supposed to be continuous by construction, and the topology is chosen to be the strongest topology with that property.

Geometrically, quotient spaces should be imagined as spaces obtained by identifying points. A circle can be formed by identifying the two endpoints of an interval. A torus can be formed by identifying opposite edges of a square. The quotient topology is the formal rule that makes these gluing pictures mathematically precise.

\begin{definition}[Quotient Topology]
Let $X$ be a topological space, let $Y$ be a set, and let $q:X\to Y$ be a surjective map. The \textbf{quotient topology} on $Y$ is defined by declaring $U\subseteq Y$ open if and only if $q^{-1}(U)$ is open in $X$.
\end{definition}

\begin{example}[Circle as a Quotient]
The circle $S^1$ can be obtained by identifying the endpoints of the interval $[0,1]$. Formally, define an equivalence relation on $[0,1]$ by $0\sim 1$ and $x\sim x$ otherwise. The quotient space $[0,1]/\sim$ is homeomorphic to $S^1$.
\end{example}

\begin{remark}
Quotient topology is the rigorous language of gluing. Many geometric spaces are built by taking a simple space and identifying selected points or subspaces.
\end{remark}

\section{Separation Axioms}


Separation axioms measure how well a topology can distinguish points and closed sets. In a Hausdorff space, two distinct points can be placed into disjoint neighborhoods. This property is so common in analysis that it is easy to overlook, but it is essential for many familiar results. For example, limits of sequences are unique in Hausdorff spaces.

Weaker topologies may fail to separate points, which means the space intentionally forgets some distinctions. Stronger separation properties make the topology behave more like metric spaces. Thus separation axioms form a scale: as we move up the scale, the topology contains enough open sets to support more refined arguments.

Topological spaces can be very strange. Separation axioms measure how well a topology distinguishes points.

\begin{definition}[$T_0$, $T_1$, and Hausdorff Spaces]
Let $X$ be a topological space.
\begin{enumerate}
    \item $X$ is $T_0$ if for any distinct $x,y\in X$, there exists an open set containing one of them but not the other.
    \item $X$ is $T_1$ if for any distinct $x,y\in X$, there exists an open set containing $x$ but not $y$, and an open set containing $y$ but not $x$.
    \item $X$ is \textbf{Hausdorff} or $T_2$ if for any distinct $x,y\in X$, there exist disjoint open sets $U,V$ such that $x\in U$ and $y\in V$.
\end{enumerate}
\end{definition}

\begin{proposition}
Every Hausdorff space is $T_1$, and every $T_1$ space is $T_0$.
\end{proposition}

\begin{example}
Every metric space is Hausdorff. If $x\neq y$, let $r=d(x,y)/3$. Then $B_r(x)$ and $B_r(y)$ are disjoint open neighborhoods.
\end{example}

\begin{proposition}
A space $X$ is $T_1$ if and only if every singleton set $\{x\}$ is closed.
\end{proposition}

\begin{proof}
If $X$ is $T_1$, then for every $y\neq x$, there is an open set $U_y$ containing $y$ but not $x$. Thus
\[
    X\setminus\{x\}=\bigcup_{y\neq x}U_y
\]
is open, so $\{x\}$ is closed. Conversely, if every singleton is closed, then $X\setminus\{y\}$ is an open set containing $x$ but not $y$ whenever $x\neq y$.
\end{proof}

\begin{theorem}[Uniqueness of Limits in Hausdorff Spaces]
If $X$ is Hausdorff, then a sequence in $X$ has at most one limit.
\end{theorem}

\begin{proof}
Suppose $x_n\to x$ and $x_n\to y$ with $x\neq y$. Since $X$ is Hausdorff, choose disjoint open sets $U$ and $V$ with $x\in U$ and $y\in V$. Eventually $x_n\in U$, and eventually $x_n\in V$. Hence eventually $x_n\in U\cap V$, impossible because $U\cap V=\emptyset$.
\end{proof}

\section{Continuity and Homeomorphism}


The topological definition of continuity is inverse-image based: $f:X\to Y$ is continuous if the preimage of every open set in $Y$ is open in $X$. This definition may look less intuitive than the $\varepsilon$--$\delta$ definition, but it captures the same idea without numerical distances. It says that whenever the target imposes a local test of nearness, the source passes that test through preimages.

Homeomorphisms are the isomorphisms of topology. If two spaces are homeomorphic, they may look geometrically different, but they have the same topological structure. A coffee mug and a torus are often used as a visual metaphor because each has one hole; topology regards them as equivalent under continuous deformation, provided tearing and gluing are not allowed.

\subsection{Continuous Functions}

\begin{definition}[Continuity]
Let $X$ and $Y$ be topological spaces. A function $f:X\to Y$ is \textbf{continuous} if for every open set $V\subseteq Y$, the preimage $f^{-1}(V)$ is open in $X$.
\end{definition}

\begin{proposition}[Equivalent Form Using Closed Sets]
A function $f:X\to Y$ is continuous if and only if for every closed set $C\subseteq Y$, the preimage $f^{-1}(C)$ is closed in $X$.
\end{proposition}

\begin{proof}
This follows from
\[
    f^{-1}(Y\setminus C)=X\setminus f^{-1}(C).
\]
\end{proof}

\begin{proposition}[Composition of Continuous Functions]
If $f:X\to Y$ and $g:Y\to Z$ are continuous, then $g\circ f:X\to Z$ is continuous.
\end{proposition}

\begin{proof}
For any open set $W\subseteq Z$,
\[
    (g\circ f)^{-1}(W)=f^{-1}(g^{-1}(W)).
\]
Since $g$ is continuous, $g^{-1}(W)$ is open in $Y$. Since $f$ is continuous, $f^{-1}(g^{-1}(W))$ is open in $X$.
\end{proof}

\subsection{Homeomorphisms}

\begin{definition}[Homeomorphism]
A function $f:X\to Y$ is a \textbf{homeomorphism} if:
\begin{enumerate}
    \item $f$ is bijective;
    \item $f$ is continuous;
    \item $f^{-1}$ is continuous.
\end{enumerate}
If such a map exists, $X$ and $Y$ are called \textbf{homeomorphic}.
\end{definition}

Homeomorphic spaces are considered the same from the topological viewpoint. They may look geometrically different, but their open-set structures are equivalent.

\begin{example}
The open intervals $(0,1)$ and $\mathbb{R}$ are homeomorphic. One possible homeomorphism is
\[
    f(x)=\tan\left(\pi x-\frac{\pi}{2}\right).
\]
\end{example}

\begin{remark}
A bijective continuous function need not be a homeomorphism. The inverse must also be continuous. However, a continuous bijection from a compact space to a Hausdorff space is always a homeomorphism.
\end{remark}

\section{Connectedness}


Connectedness formalizes the idea of being in one piece. A space is disconnected if it can be separated into two non-empty open parts that do not interact. The definition is deliberately phrased in terms of open sets rather than pictures, because topological spaces need not come with an obvious visual representation.

The intermediate value theorem is a fundamental example of connectedness in action. The interval $[a,b]$ is connected, and a continuous image of a connected set is connected. Since connected subsets of $\mathbb{R}$ are intervals, a continuous real-valued function on $[a,b]$ cannot jump over intermediate values. Thus a classical theorem of calculus becomes a theorem about preservation of connectedness.

Connectedness formalizes the idea that a space consists of one piece.

\subsection{Definitions}

\begin{definition}[Separation]
A \textbf{separation} of a topological space $X$ is a pair of non-empty disjoint open sets $U,V\subseteq X$ such that
\[
    X=U\cup V.
\]
\end{definition}

\begin{definition}[Connected Space]
A topological space $X$ is \textbf{connected} if it has no separation.
\end{definition}

\begin{proposition}[Equivalent Characterization]
A space $X$ is connected if and only if the only subsets of $X$ that are both open and closed are $\emptyset$ and $X$.
\end{proposition}

\begin{proof}
If $A\subseteq X$ is both open and closed and is neither $\emptyset$ nor $X$, then $A$ and $X\setminus A$ form a separation. Conversely, if $X=U\cup V$ is a separation, then $U$ is open and also closed because $U=X\setminus V$.
\end{proof}

\begin{theorem}[Intervals Are Connected]
Every interval in $\mathbb{R}$ is connected.
\end{theorem}

\begin{proof}
We prove the idea for an interval $I$. Suppose $I=A\cup B$ is a separation, with $a\in A$ and $b\in B$, say $a<b$. Let
\[
    S=A\cap [a,b].
\]
Using the least upper bound property of $\mathbb{R}$, one shows that $c=\sup S$ must belong neither consistently to $A$ nor to $B$ without contradicting the openness of $A$ and $B$ in the subspace $I$. Hence no such separation exists.
\end{proof}

\subsection{Continuous Images of Connected Spaces}

\begin{theorem}
If $f:X\to Y$ is continuous and $X$ is connected, then $f(X)$ is connected.
\end{theorem}

\begin{proof}
Suppose $f(X)$ were disconnected. Then $f(X)=U\cup V$ for non-empty disjoint open sets in the subspace topology of $f(X)$. The preimages $f^{-1}(U)$ and $f^{-1}(V)$ are non-empty disjoint open subsets of $X$ whose union is $X$, contradicting connectedness.
\end{proof}

\begin{theorem}[Intermediate Value Theorem]
Let $f:[a,b]\to\mathbb{R}$ be continuous. If $f(a)<y<f(b)$ or $f(b)<y<f(a)$, then there exists $c\in(a,b)$ such that $f(c)=y$.
\end{theorem}

\begin{proof}
The interval $[a,b]$ is connected, and the continuous image $f([a,b])$ is connected in $\mathbb{R}$. Connected subsets of $\mathbb{R}$ are intervals. Hence $f([a,b])$ contains every value between $f(a)$ and $f(b)$.
\end{proof}

\subsection{Path Connectedness}

\begin{definition}[Path Connectedness]
A space $X$ is \textbf{path connected} if for any $x,y\in X$, there exists a continuous map $\gamma:[0,1]\to X$ such that $\gamma(0)=x$ and $\gamma(1)=y$.
\end{definition}

\begin{proposition}
Every path connected space is connected.
\end{proposition}

\begin{proof}
Suppose $X$ is path connected but disconnected, with separation $X=U\cup V$. Choose $x\in U$ and $y\in V$. Let $\gamma:[0,1]\to X$ be a path from $x$ to $y$. Then $\gamma([0,1])$ is connected as a continuous image of a connected interval, but it intersects both $U$ and $V$, giving a separation of $\gamma([0,1])$. This is impossible.
\end{proof}

\begin{remark}
The converse is false. There exist connected spaces that are not path connected, such as the topologist's sine curve. Thus path connectedness is stronger than connectedness.
\end{remark}

\section{Compactness}


Compactness is one of the most important replacement principles in analysis. In $\mathbb{R}^n$, compactness is equivalent to being closed and bounded, but the open-cover definition is more fundamental. It says that whenever a space is covered by open sets, finitely many of those sets already suffice. Compactness is therefore a finiteness principle hidden inside an infinite setting.

Many theorems in analysis rely on this finite extraction property. Continuous functions on compact spaces attain maxima and minima because open neighborhoods that control local behavior can be reduced to finitely many neighborhoods. Uniform continuity on compact metric spaces follows for the same reason: local control can be upgraded to global control by compactness.

Compactness is one of the central ideas of topology and analysis. In finite sets, many arguments work because a finite cover has finitely many pieces. Compactness generalizes this finite behavior to infinite spaces.

\subsection{Open Covers}

\begin{definition}[Open Cover]
Let $X$ be a topological space and $A\subseteq X$. An \textbf{open cover} of $A$ is a collection of open sets $\{U_i\}_{i\in I}$ such that
\[
    A\subseteq\bigcup_{i\in I}U_i.
\]
A \textbf{subcover} is a subcollection that still covers $A$.
\end{definition}

\begin{definition}[Compactness]
A topological space $X$ is \textbf{compact} if every open cover of $X$ has a finite subcover.
\end{definition}

\begin{example}
Every finite topological space is compact. Given any open cover, choose one open set containing each point. Since there are finitely many points, this gives a finite subcover.
\end{example}

\begin{example}
The open interval $(0,1)$ is not compact in the usual topology. The collection
\[
    \left\{\left(\frac{1}{n},1\right):n\geq 2\right\}
\]
is an open cover of $(0,1)$, but no finite subcollection covers points sufficiently close to $0$.
\end{example}

\subsection{Compactness in $\mathbb{R}^n$}

\begin{theorem}[Heine-Borel Theorem]
A subset $K\subseteq\mathbb{R}^n$ is compact if and only if it is closed and bounded.
\end{theorem}

\begin{remark}
The Heine-Borel theorem is special to finite-dimensional Euclidean spaces. In general metric spaces, closed and bounded does not necessarily imply compact.
\end{remark}

\begin{theorem}[Bolzano-Weierstrass Theorem]
Every bounded sequence in $\mathbb{R}^n$ has a convergent subsequence.
\end{theorem}

\begin{remark}
In $\mathbb{R}^n$, compactness, sequential compactness, and closed boundedness are equivalent. In arbitrary topological spaces, these equivalences may fail.
\end{remark}

\subsection{Continuous Images of Compact Spaces}

\begin{theorem}
If $f:X\to Y$ is continuous and $X$ is compact, then $f(X)$ is compact.
\end{theorem}

\begin{proof}
Let $\{V_i\}_{i\in I}$ be an open cover of $f(X)$. Then $\{f^{-1}(V_i)\}_{i\in I}$ is an open cover of $X$. Since $X$ is compact, finitely many preimages cover $X$, say $f^{-1}(V_{i_1}),\dots,f^{-1}(V_{i_n})$. Then $V_{i_1},\dots,V_{i_n}$ cover $f(X)$.
\end{proof}

\begin{theorem}[Extreme Value Theorem]
Let $X$ be compact and let $f:X\to\mathbb{R}$ be continuous. Then $f$ attains a maximum and a minimum on $X$.
\end{theorem}

\begin{proof}
The image $f(X)$ is compact in $\mathbb{R}$. By Heine-Borel, it is closed and bounded. Since it is bounded, it has a supremum and an infimum. Since it is closed, these values belong to $f(X)$. Therefore $f$ attains its maximum and minimum.
\end{proof}

\begin{theorem}[Compact to Hausdorff]
Let $X$ be compact, $Y$ Hausdorff, and $f:X\to Y$ a continuous bijection. Then $f$ is a homeomorphism.
\end{theorem}

\begin{proof}
It suffices to show that $f$ is a closed map. Let $C\subseteq X$ be closed. Since $X$ is compact, $C$ is compact. The continuous image $f(C)$ is compact in $Y$. In a Hausdorff space, compact subsets are closed. Hence $f(C)$ is closed, so $f^{-1}$ is continuous.
\end{proof}

\section{Metric Spaces Revisited}


Returning to metric spaces after developing topology is useful because it clarifies what extra structure a metric provides. A metric gives numerical distances, sequences, Cauchy conditions, completeness, and quantitative estimates. A topology gives open sets and continuity. Every metric determines a topology, but not every topology comes from a metric.

This distinction matters in advanced mathematics. Some spaces are naturally topological but not naturally metric; others admit many different metrics producing the same topology. When a proof uses only open sets, it is topological. When it uses distances, Cauchy sequences, or explicit bounds, it is metric. Recognizing this difference helps identify the true assumptions behind a theorem.

Many topological spaces come from metrics, but not all topologies are metrizable.

\begin{definition}[Metrizable Space]
A topological space $(X,\mathcal{T})$ is \textbf{metrizable} if there exists a metric $d$ on $X$ such that $\mathcal{T}$ is exactly the topology generated by the open balls of $d$.
\end{definition}

\begin{example}
Every discrete topological space is metrizable. Define
\[
    d(x,y)=
    \begin{cases}
    0,&x=y,\\
    1,&x\neq y.
    \end{cases}
\]
Then every singleton is open, so every subset is open.
\end{example}

\begin{remark}
Metrizability is a strong property. Metric spaces are always Hausdorff and first countable. Therefore any topology failing these properties cannot come from a metric.
\end{remark}


\section{Separation Axioms Beyond the Hausdorff Condition}

The axioms $T_0$, $T_1$, and $T_2$ compare neighborhoods of points. The next separation axioms compare points with closed sets, and then compare two closed sets with each other. This progression is not cosmetic. Many constructions in analysis require continuous functions that distinguish sets, and such functions cannot be produced from the Hausdorff condition alone.

There are several conventions in the literature. In these notes, a \emph{regular space} and a \emph{normal space} are not assumed to be $T_1$ unless this is stated explicitly. Thus $T_3$ means regular plus $T_1$, and $T_4$ means normal plus $T_1$.

\begin{definition}[Regularity and Normality]
Let $X$ be a topological space.
\begin{enumerate}
    \item $X$ is \textbf{regular} if whenever $x\in X$ and $F\subseteq X$ is closed with $x\notin F$, there exist disjoint open sets $U,V$ such that $x\in U$ and $F\subseteq V$.
    \item $X$ is \textbf{completely regular} if whenever $x\notin F$ and $F$ is closed, there exists a continuous function $f:X\to[0,1]$ such that $f(x)=0$ and $f|_F=1$.
    \item $X$ is \textbf{normal} if whenever $A,B\subseteq X$ are disjoint closed sets, there exist disjoint open sets $U,V$ such that $A\subseteq U$ and $B\subseteq V$.
\end{enumerate}
\end{definition}

\begin{definition}[The Separation Hierarchy]
A space is called
\[
\begin{array}{c|l}
T_0 & \text{Kolmogorov},\\
T_1 & \text{all singleton sets are closed},\\
T_2 & \text{Hausdorff},\\
T_3 & \text{regular and }T_1,\\
T_{3\frac12} & \text{completely regular and }T_1,\\
T_4 & \text{normal and }T_1.
\end{array}
\]
A $T_{3\frac12}$ space is also called a \textbf{Tychonoff space}.
\end{definition}

\begin{proposition}[Closure Form of Regularity]
A space $X$ is regular if and only if for every point $x$ and every open set $U$ containing $x$, there exists an open set $V$ such that
\[
    x\in V\subseteq \overline{V}\subseteq U.
\]
\end{proposition}

\begin{proof}
Assume first that $X$ is regular. The set $F=X\setminus U$ is closed and does not contain $x$. Choose disjoint open sets $V,W$ with $x\in V$ and $F\subseteq W$. Since $V\cap W=\emptyset$, no point of $W$ lies in $\overline{V}$. Therefore $\overline{V}\subseteq X\setminus W\subseteq U$.

Conversely, suppose the closure condition holds. Let $F$ be closed and $x\notin F$. Apply the condition to $U=X\setminus F$ and choose $V$ with $x\in V\subseteq\overline{V}\subseteq U$. Then $V$ and $X\setminus\overline{V}$ are disjoint open sets containing $x$ and $F$, respectively.
\end{proof}

\begin{theorem}[Metric Spaces Are Normal]
Every metric space is normal. Consequently, every metric space is $T_4$ and hence also $T_{3\frac12}$, $T_3$, $T_2$, $T_1$, and $T_0$.
\end{theorem}

\begin{proof}
Let $A$ and $B$ be disjoint closed subsets of a metric space $(X,d)$. Define
\[
    d(x,A)=\inf_{a\in A}d(x,a),\qquad d(x,B)=\inf_{b\in B}d(x,b).
\]
Both functions are continuous. Set
\[
    U=\{x:d(x,A)<d(x,B)\},\qquad
    V=\{x:d(x,B)<d(x,A)\}.
\]
These sets are open and disjoint. If $a\in A$, then $d(a,A)=0$, while $d(a,B)>0$ because $B$ is closed and $a\notin B$. Thus $A\subseteq U$, and similarly $B\subseteq V$.
\end{proof}

\begin{example}[Why the Hypotheses Matter]
The cofinite topology on an infinite set is $T_1$ but not Hausdorff: any two non-empty open sets intersect. Thus the implication $T_1\Rightarrow T_2$ fails. The Sierpinski space $X=\{0,1\}$ with topology $\{\emptyset,\{1\},X\}$ is $T_0$ but not $T_1$. These examples show that the lower separation axioms are genuinely distinct.
\end{example}

\begin{theorem}[Urysohn's Lemma]
Let $X$ be a normal space and let $A,B\subseteq X$ be disjoint closed sets. Then there exists a continuous function $f:X\to[0,1]$ such that
\[
    f|_A=0,\qquad f|_B=1.
\]
\end{theorem}

\begin{proof}[Construction Outline]
The proof repeatedly uses normality to insert closures between open sets. For each dyadic rational $r\in[0,1]$, construct an open set $U_r$ so that
\[
    A\subseteq U_0,
    \qquad \overline{U_r}\subseteq U_s\quad(r<s),
    \qquad X\setminus U_1\supseteq B.
\]
Define
\[
    f(x)=\inf\{r:x\in U_r\},
\]
with the convention that the infimum of the empty set is $1$. The nesting relation implies
\[
    f^{-1}([0,a))=\bigcup_{r<a}U_r,
    \qquad
    f^{-1}((a,1])=\bigcup_{r>a}(X\setminus\overline{U_r}),
\]
so both relevant inverse images are open. Hence $f$ is continuous and has the required boundary values.
\end{proof}

\begin{theorem}[Tietze Extension Theorem]
Let $X$ be a normal space, let $A\subseteq X$ be closed, and let $f:A\to[-1,1]$ be continuous. Then there exists a continuous function $F:X\to[-1,1]$ such that $F|_A=f$.
\end{theorem}

\begin{remark}
Urysohn's lemma separates two closed sets by a continuous real-valued function. Tietze's theorem extends an already defined continuous function from a closed subspace. The second theorem is deeper, but its proof repeatedly applies the first theorem to approximate the remaining error. These results explain why normality is the natural topological condition for constructing real-valued functions.
\end{remark}

\section{Components, Local Connectedness, and the Topologist's Sine Curve}

Connectedness is a global condition: the whole space cannot be separated into two non-empty open pieces. Path connectedness is more constructive: every pair of points can be joined by a continuous path. Local connectedness asks for connected neighborhoods near each point. None of these conditions should be confused with the others.

\begin{definition}[Connected and Path Components]
For $x\in X$, the \textbf{connected component} of $x$ is the union of all connected subsets of $X$ containing $x$. The \textbf{path component} of $x$ is the union of all path connected subsets containing $x$.
\end{definition}

\begin{proposition}
Connected components are closed, pairwise disjoint, and maximal connected subsets. Path components are pairwise disjoint and maximal path connected subsets, but they need not be closed in an arbitrary topological space.
\end{proposition}

\begin{proof}
If two connected subsets intersect, their union is connected. Hence the union used to define a connected component is connected, and two different components cannot intersect. The closure of a connected set is connected, so maximality forces each connected component to contain its closure and therefore to be closed. The path-component statement follows from the analogous fact that two intersecting path connected sets have path connected union.
\end{proof}

\begin{definition}[Local Connectedness]
A space $X$ is \textbf{locally connected at $x$} if every neighborhood of $x$ contains a connected open neighborhood of $x$. It is \textbf{locally path connected at $x$} if every neighborhood of $x$ contains a path connected open neighborhood of $x$. The space is locally connected or locally path connected if the condition holds at every point.
\end{definition}

\begin{theorem}
If $X$ is locally path connected, then its connected components and path components coincide, and each component is open.
\end{theorem}

\begin{proof}
Fix $x\in X$ and let $P$ be its path component. For each $y\in P$, choose a path connected open neighborhood $U_y$ of $y$. Every point of $U_y$ can be joined to $y$, and $y$ can be joined to $x$, so $U_y\subseteq P$. Hence $P$ is open. The path components form a disjoint open partition. A connected subset cannot meet two different members of such a partition, so the connected component of $x$ is contained in $P$. The reverse containment is automatic because path connected sets are connected.
\end{proof}

\begin{example}[The Topologist's Sine Curve]
Let
\[
    S=\left\{\left(x,\sin\frac1x\right):0<x\leq 1\right\}
      \cup\left(\{0\}\times[-1,1]\right)\subseteq\mathbb{R}^2.
\]
The graph part
\[
    G=\left\{\left(x,\sin\frac1x\right):0<x\leq 1\right\}
\]
is the continuous image of the connected interval $(0,1]$, so $G$ is connected. Its closure is exactly $S$, and the closure of a connected set is connected. Therefore $S$ is connected.

However, no path in $S$ joins a point of the vertical segment to a point of $G$. Suppose such a path existed. As the first coordinate approaches $0$ through positive values, the second coordinate would have to follow $\sin(1/x)$, which oscillates infinitely often between values near $-1$ and $1$. Continuity at the first moment when the path leaves the vertical segment would be impossible. Hence $S$ is connected but not path connected. It is also not locally connected at points of the vertical segment.
\end{example}

\begin{remark}
This example is valuable because the failure is not caused by disconnected pieces placed far apart. The graph accumulates on the entire vertical segment, so closure creates connectedness without creating paths. It is a model example of why limit behavior in topology can be subtler than geometric intuition suggests.
\end{remark}

\section{Compactness Beyond Closed and Bounded Sets}

The slogan ``compact means closed and bounded'' is valid in finite-dimensional Euclidean space, but it is not the definition and it fails in many other spaces. The open-cover definition is robust because it is preserved by continuous maps and makes sense without a metric. Several equivalent formulations reveal different ways in which compactness converts infinitely many local conditions into finitely many global conditions.

\subsection{The Finite Intersection Property}

\begin{definition}[Finite Intersection Property]
A family $\mathcal{F}$ of subsets of $X$ has the \textbf{finite intersection property} if every finite subfamily has non-empty intersection.
\end{definition}

\begin{theorem}[Closed-Set Characterization of Compactness]
A topological space $X$ is compact if and only if every family of closed sets with the finite intersection property has non-empty total intersection.
\end{theorem}

\begin{proof}
Suppose $X$ is compact and $\{F_i\}_{i\in I}$ is a family of closed sets with empty total intersection. Then $\{X\setminus F_i\}_{i\in I}$ is an open cover of $X$. A finite subcover would imply that the corresponding finite intersection of the $F_i$ is empty, contradicting the finite intersection property.

Conversely, suppose every closed family with the finite intersection property has non-empty intersection. If an open cover $\{U_i\}_{i\in I}$ had no finite subcover, then every finite intersection of the closed sets $X\setminus U_i$ would be non-empty. Their total intersection would therefore be non-empty, contradicting the assumption that the $U_i$ cover $X$.
\end{proof}

\subsection{Sequential Compactness in Metric Spaces}

\begin{definition}[Sequential Compactness]
A space is \textbf{sequentially compact} if every sequence has a convergent subsequence whose limit belongs to the space.
\end{definition}

\begin{theorem}
For metric spaces, compactness and sequential compactness are equivalent.
\end{theorem}

\begin{proof}[Idea of Proof]
If a metric space is compact, cover it by finitely many balls of radius $1$. One of these balls contains infinitely many terms of the sequence. Cover that ball by finitely many balls of radius $1/2$ and choose one containing infinitely many of the remaining terms. Continuing inductively gives a nested selection of terms whose mutual distances tend to zero. Compactness supplies a cluster point, and the selected subsequence converges to it.

For the converse, assume sequential compactness. First, the space is totally bounded. Otherwise there would exist $\varepsilon>0$ and a sequence $(x_n)$ with
\[
    d(x_m,x_n)\geq\varepsilon\qquad(m\neq n),
\]
which has no convergent subsequence.

Now let $\mathcal{U}$ be an open cover. We claim that $\mathcal{U}$ has a Lebesgue number. If not, for each $n$ choose $x_n$ such that $B_{1/n}(x_n)$ is contained in no member of $\mathcal{U}$. Sequential compactness gives a subsequence $x_{n_k}\to x$. Choose $U\in\mathcal{U}$ and $r>0$ with $B_r(x)\subseteq U$. For large $k$, both $d(x_{n_k},x)<r/2$ and $1/n_k<r/2$, so
\[
    B_{1/n_k}(x_{n_k})\subseteq B_r(x)\subseteq U,
\]
a contradiction. Finally, total boundedness supplies finitely many balls of radius smaller than the Lebesgue number, and one member of $\mathcal{U}$ contains each ball. These finitely many members form a finite subcover.
\end{proof}

\begin{theorem}[Lebesgue Number Lemma]
Let $(X,d)$ be a compact metric space and let $\mathcal{U}$ be an open cover of $X$. Then there exists $\delta>0$ such that every subset of diameter less than $\delta$ is contained in some member of $\mathcal{U}$.
\end{theorem}

\begin{proof}
For each $x\in X$, choose $U_x\in\mathcal{U}$ containing $x$, and choose $r_x>0$ such that $B_{2r_x}(x)\subseteq U_x$. The balls $B_{r_x}(x)$ cover $X$, so finitely many suffice: $B_{r_1}(x_1),\dots,B_{r_n}(x_n)$. Let $\delta=\min r_i$. If a set $A$ has diameter less than $\delta$ and contains a point $a$, choose $i$ with $a\in B_{r_i}(x_i)$. Then every $y\in A$ satisfies
\[
    d(y,x_i)\leq d(y,a)+d(a,x_i)<\delta+r_i\leq 2r_i,
\]
so $A\subseteq B_{2r_i}(x_i)\subseteq U_{x_i}$.
\end{proof}

\begin{corollary}[Uniform Continuity on Compact Spaces]
Every continuous map from a compact metric space to a metric space is uniformly continuous.
\end{corollary}

\begin{proof}
Let $f:X\to Y$ be continuous and fix $\varepsilon>0$. For each $x\in X$, choose an open neighborhood $U_x$ such that $f(U_x)$ lies in the ball of radius $\varepsilon/2$ around $f(x)$. A Lebesgue number $\delta$ for the cover $\{U_x\}$ ensures that whenever $d(p,q)<\delta$, the two-point set $\{p,q\}$ lies in one $U_x$. Hence
\[
    d_Y(f(p),f(q))<\varepsilon.
\]
\end{proof}

\subsection{Finite Products and the First Form of Tychonoff's Theorem}

\begin{theorem}[Finite Product Theorem]
If $X$ and $Y$ are compact, then $X\times Y$ is compact in the product topology. Consequently, every finite product of compact spaces is compact.
\end{theorem}

\begin{proof}[Tube-Lemma Proof]
Let $\mathcal{U}$ be an open cover of $X\times Y$. Fix $x\in X$. The slice $\{x\}\times Y$ is homeomorphic to $Y$ and hence compact, so finitely many members of $\mathcal{U}$ cover it. Their union is an open set containing $\{x\}\times Y$. By the tube lemma, there exists an open neighborhood $N_x$ of $x$ such that
\[
    N_x\times Y
\]
is contained in that finite union. The sets $N_x$ cover $X$. Compactness of $X$ gives finitely many of them, and collecting the corresponding finite subfamilies of $\mathcal{U}$ yields a finite cover of $X\times Y$.
\end{proof}

\begin{lemma}[Tube Lemma]
Let $Y$ be compact and let $W\subseteq X\times Y$ be open with $\{x_0\}\times Y\subseteq W$. Then there exists an open neighborhood $N$ of $x_0$ such that $N\times Y\subseteq W$.
\end{lemma}

\begin{proof}
For each $y\in Y$, choose basic open sets $N_y\subseteq X$ and $V_y\subseteq Y$ such that
\[
    (x_0,y)\in N_y\times V_y\subseteq W.
\]
The sets $V_y$ cover $Y$, so finitely many $V_{y_1},\dots,V_{y_n}$ suffice. Then
\[
    N=N_{y_1}\cap\cdots\cap N_{y_n}
\]
is an open neighborhood of $x_0$, and $N\times Y\subseteq W$.
\end{proof}

\begin{remark}
The full Tychonoff theorem states that an arbitrary product of compact spaces is compact. Its standard proof uses the axiom of choice, often through ultrafilters or maximal families with the finite intersection property. The finite theorem above contains the geometric core and is sufficient for most elementary applications.
\end{remark}

\section{Products and Quotients as Universal Constructions}

The definitions of product and quotient topologies can appear asymmetric. The product topology is generated by inverse images under projections, while the quotient topology is tested by inverse images under one surjective map. Their common feature is a universal property: each topology is exactly the one that makes a certain continuity problem as easy as possible.

\begin{theorem}[Universal Property of Products]
Let $Z$ be a topological space and let $f:Z\to X\times Y$. Write
\[
    f(z)=(f_1(z),f_2(z)).
\]
Then $f$ is continuous if and only if both coordinate maps $f_1:Z\to X$ and $f_2:Z\to Y$ are continuous.
\end{theorem}

\begin{proof}
If $f$ is continuous, then $f_1=\pi_X\circ f$ and $f_2=\pi_Y\circ f$ are continuous. Conversely, for every basic open set $U\times V$,
\[
    f^{-1}(U\times V)=f_1^{-1}(U)\cap f_2^{-1}(V),
\]
which is open when $f_1$ and $f_2$ are continuous.
\end{proof}

\begin{theorem}[Universal Property of Quotients]
Let $q:X\to Y$ be a quotient map and let $g:Y\to Z$ be a function. Then $g$ is continuous if and only if $g\circ q:X\to Z$ is continuous.
\end{theorem}

\begin{proof}
One direction follows from composition. Conversely, assume $g\circ q$ is continuous. For every open set $W\subseteq Z$,
\[
    q^{-1}(g^{-1}(W))=(g\circ q)^{-1}(W)
\]
is open in $X$. By the definition of the quotient topology, $g^{-1}(W)$ is open in $Y$.
\end{proof}

\begin{example}[The Circle from an Interval]
Define $q:[0,1]\to S^1$ by
\[
    q(t)=e^{2\pi i t}.
\]
The map is continuous, surjective, and identifies exactly the endpoints $0$ and $1$. Since $[0,1]$ is compact and $S^1$ is Hausdorff, the induced bijection
\[
    [0,1]/(0\sim1)\longrightarrow S^1
\]
is a homeomorphism. This argument is stronger than a picture: it proves that the quotient topology is precisely the usual topology on the circle.
\end{example}

\begin{example}[A Non-Hausdorff Quotient]
Start with two copies of the real line and identify corresponding points except at the origins. The resulting ``line with two origins'' is locally indistinguishable from $\mathbb{R}$, but the two origins cannot be separated by disjoint open sets. Quotients can therefore destroy the Hausdorff property. Whenever a space is built by gluing, separation axioms must be checked rather than assumed.
\end{example}

\section{Conclusion}

Topology begins by replacing distance with open sets. From this simple abstraction, we recover the language of continuity, convergence, compactness, and connectedness. The power of topology is that it reveals which theorems of analysis depend on metric quantities and which depend only on the qualitative structure of open sets.

The essential pattern is the same as in previous chapters: identify the primitive data, state axioms, define structure-preserving maps, and then study properties invariant under those maps. In topology, the primitive data are open sets, the structure-preserving maps are continuous functions, and the notion of sameness is homeomorphism.
